\section{Discussion and Limitations}
\label{sec:discussion}

\paragraph{Where the annotation saving comes from.}
The system does not make geometry disappear; it changes its amortization.  A
manual workflow calibrates each image or clicks its points.  Our workflow pays
for video acquisition, reconstruction, and one audited scene alignment, after
which camera placement, pose, intrinsics, point identity, line/segmentation
targets, depth, geometric validity, and renderer support are computed for every
accepted render.  The
saving grows with the number and diversity of rendered views.  This makes
partial, elevated, off-axis, and multi-court views---precisely the views useful
for defeating local shortcuts---cheap to add.

\paragraph{3DGS is necessary but not sufficient.}
Novel-view rendering alone supplies neither metres nor court semantics.  A
beautiful but unaligned Gaussian scene cannot supervise translation or physical
keypoint identity.  Conversely, a court transform without a renderer produces
only labels, not realistic pixels at novel poses.  The contribution lies in
their composition through one transform graph and in rejecting the composition
when scale, topology, or held-out evidence is weak.

\paragraph{Alignment bootstrapping and reconstruction failure.}
Alignment consumes a learned line map, so its error modes include those of the
existing line detector.  Sparse reconstruction points may poorly support the
ground, reflective/occluded lines may be missed, and repetitive neighbouring
courts can induce 90-degree or court-identity mistakes.  The common-scale,
whole-template, topology, and holdout gates reduce propagation of these errors
but also reject useful-looking scenes.  They are quality controls, not a proof
that every accepted transform is exact.

\paragraph{Rendering domain gap.}
A 3DGS reconstruction stays close to the captured venue but may contain floaters,
blur, missing transient objects, or unreliable extrapolation far outside the
captured camera envelope.  Trajectory radii are therefore anchored to captured
offsets and geometry is screened before rendering.  Training-time photometric
augmentation provides additional variation, but this paper does not claim that
it closes the gap to every broadcast camera, surface, weather condition, or
compression pipeline.

\paragraph{Camera model.}
The pose head assumes equal focal lengths and uses an augmentation-provided
principal point.  It does not predict distortion, skew, or a moving principal
point.  Unsupported transforms are rejected so that the labels remain unique,
but broader camera models require an expanded head and projection loss.  The
camera-end canonicalization also assumes the two proper choices
\(I/R_z(\pi)\); severely ambiguous mid-plane views must be filtered rather than
assigned silently.

\paragraph{Scope of shortcut resistance.}
The global query and reprojection objective constrain the solution class, but
they do not prove that a neural network ignores every spurious cue.  The network
may exploit renderer artifacts or venue-specific texture correlated with pose,
and the two branches may co-adapt.  Direct ground-truth losses, independent
pose-reprojection metrics, stop-gradient ablations, worst-stratum reporting,
and the local-context diagnostic are necessary to test the intended mechanism.
Held-out real calibrated data remain the strongest missing evaluation.

\paragraph{Compute and data governance.}
Reconstruction and novel-view rendering add cost before training, although one
scene can amortize that cost over many task datasets and camera views.  Source
videos and reconstructed appearance may contain people or venue-identifying
information; dataset publication must respect the source license and privacy
constraints independently of the software's open-source license.
